\section{Conclusion}
\label{sec:conclusion}

We asked when a language model should revise its own answer. Our analytical model shows that confidence gating can only help when the revision step actually repairs errors and the confidence signal is informative in the relevant region of its ROC curve. An uninformative signal can never beat the better of never and always revising, and the achievable gain is capped by the fix rate. Experiments with \modelA{} confirm both sides of this picture. Self-critique almost never repairs an error, so no gate improves on never revising, even though self-consistency identifies wrong answers well. When the majority answer of the self-consistency samples is used as the revision step instead, the same gates behave as the model predicts. For this reason, future work should focus on the revision step itself, and on selecting thresholds reliably from small development sets.
